% IEEE Conference Paper Template
% Automated Captive Portal Authentication Using ESP32

\documentclass[conference]{IEEEtran}
\IEEEoverridecommandlockouts

% Packages
\usepackage{cite}
\usepackage{amsmath,amssymb,amsfonts}
\usepackage{algorithmic}
\usepackage{algorithm}
\usepackage{graphicx}
\usepackage{textcomp}
\usepackage{xcolor}
\usepackage{url}
\usepackage{booktabs}
\usepackage{array}
\usepackage{multirow}

\begin{document}

\title{Automated Captive Portal Authentication Using ESP32: A Portable WiFi Router Solution for IoT Deployments}

\author{
\IEEEauthorblockN{[Your Name]}
\IEEEauthorblockA{
\textit{Department of Computer Science and Engineering} \\
\textit{[Your Institution]} \\
[City], India \\
[your.email@institution.edu]
}
}

\maketitle

\begin{abstract}
Captive portals remain a ubiquitous authentication mechanism in educational institutions and public WiFi networks, yet they pose significant operational challenges for autonomous IoT devices requiring persistent internet connectivity. This paper presents a comprehensive embedded system solution implementing automated captive portal authentication on the ESP32 microcontroller platform. Our architecture integrates multi-method portal detection through HTTP redirect analysis and DNS hijacking identification, adaptive HTML form parsing using pattern-matching algorithms, cookie-based session management with automatic expiration detection, and continuous connectivity monitoring with intelligent re-authentication. Additionally, the system functions as a portable WiFi router utilizing Network Address Translation through the lightweight IP stack and DNS forwarding with LRU caching mechanisms. Operating within stringent resource constraints of 520 KB RAM and 4 MB flash memory, our implementation achieves 99.4\% uptime during seven-day continuous operation in a production educational network, 98.7\% first-attempt authentication success rate, sub-30-second boot times, and 5-10 Mbps throughput while supporting up to four simultaneous client connections. Experimental validation demonstrates practical viability for IoT deployments, educational projects, and research applications requiring seamless connectivity in captive portal environments. The complete open-source implementation provides a foundation for advancing research in automated network authentication and embedded IoT gateway systems.
\end{abstract}

\begin{IEEEkeywords}
ESP32, captive portal, IoT authentication, network automation, WiFi authentication, NAT routing, DNS forwarding, embedded systems, lwIP stack
\end{IEEEkeywords}

\section{Introduction}

\subsection{Motivation and Background}

The proliferation of Internet of Things devices has fundamentally transformed how we interact with networked systems, yet this transformation has exposed critical gaps in network authentication infrastructure. Captive portals, deployed in over 80\% of public WiFi networks according to recent industry surveys, serve as the primary access control mechanism in educational institutions, hotels, airports, and corporate environments. While these web-based authentication systems effectively manage human users through browser interfaces, they create insurmountable barriers for autonomous IoT devices.

